\documentclass[11pt]{article}
\usepackage[margin=1in]{geometry}
\usepackage{amsmath,amssymb}
\usepackage{booktabs,longtable,array}
\usepackage{minted}
\usepackage{xcolor}
\usepackage[hidelinks]{hyperref}
\definecolor{codegray}{gray}{0.6}


%%%%%%%%%%%%%%%%%%%%%%%%%%%%%%%%%%%%%%%%%%%%%%%%%%%%%%%%%%%%%%%%
% Code listings: minted (requires -shell-escape)
%%%%%%%%%%%%%%%%%%%%%%%%%%%%%%%%%%%%%%%%%%%%%%%%%%%%%%%%%%%%%%%%

\definecolor{LightGray}{rgb}{0.95,0.95,0.95}

\setminted{
  bgcolor=LightGray,
  bgcolorpadding=6pt,
  breaklines,
  fontsize=\footnotesize,
  xleftmargin=8pt,
  xrightmargin=8pt
}


\title{MiniOps.jl: Operator and Solver Reference}
\author{Mauricio D. Sacchi}
\date{\today}

\begin{document}
\maketitle

\section{Purpose}
MiniOps is a compact matrix-free operator package for signal processing and
inverse problems.  An operator $A$ stores a forward action and its adjoint,
so that the principal syntax is
\[
    y = A x, \qquad x_a = A^* y,
\]
implemented in Julia as \texttt{A * x} and \texttt{A' * y}.  Operators can
also be composed without explicitly forming matrices.

\section{Public API}
The table below lists the symbols exported by \texttt{MiniOps.jl}.  This is
the list to check when deciding what belongs to the supported public API.

\renewcommand{\arraystretch}{1.18}
\begin{longtable}{>{\ttfamily}p{0.26\textwidth}p{0.66\textwidth}}
\toprule
\normalfont Name & Purpose \\
\midrule
\endhead
Op & Matrix-free linear operator containing forward and adjoint actions, flattened sizes, and a name. \\
with\_shape & Infer flattened operator dimensions from a prototype input. \\
conv1d\_op & Full 1-D convolution with a fixed filter; adjoint is correlation with the conjugated filter. \\
conv1d\_cols\_op & Apply the same 1-D convolution independently to every matrix column/trace. \\
sampling\_op & Scalar/vector sampling for 1-D inputs; for arrays of dimension $\ge 2$, sample complete traces while preserving dimension 1 as time. \\
trace\_sampling\_op & Explicit whole-trace sampling for arrays shaped \texttt{(nt,n2,n3,...)}. Recommended when trace semantics should be unambiguous. \\
blending\_op & Delay and sum source traces in a common-receiver gather, with fractional delays handled by linear interpolation. \\
scaling\_op & Multiply an array by a scalar; the adjoint uses the complex conjugate scalar. \\
diag\_op & Elementwise diagonal operator; the adjoint uses conjugated weights. \\
fft\_op & N-D FFT operator, unitary by default; planned and unplanned constructors are available. \\
pad\_op & Zero-pad an N-D array; the adjoint extracts the unpadded interior. \\
nudft\_nffts & Number of FFTs required by the Taylor-series NUDFT approximation for a given dimension and order. \\
nudft1\_op & 1-D Taylor/FFT approximation to the nonuniform DFT. \\
nudft2\_op & 2-D Taylor/FFT approximation to the nonuniform DFT. \\
nudft3\_op & 3-D Taylor/FFT approximation to the nonuniform DFT. \\
nudft\_txy\_op & Time--space operator combining a spatial 2-D NUDFT with a unitary FFT/IFFT in time. \\
radon\_tx\_parab\_op & Parabolic Radon transform in $(t,x)$ using linear time interpolation. \\
radon\_tx\_tp\_chirp\_op & Frequency-domain linear Radon transform using chirp factorization. \\
adjoint\_test & Numerical check of $\langle Ax,y\rangle=\langle x,A^*y\rangle$. \\
linearity\_test & Numerical check of operator linearity. \\
opnorm\_power & Power-iteration estimate of $\|A\|_2$. \\
is\_selfadjoint & Numerical comparison of $Ax$ and $A^*x$. \\
soft\_threshold & Elementwise soft thresholding. \\
hard\_threshold & Elementwise hard thresholding. \\
ista & Iterative soft-thresholding solver for an $\ell_1$-regularized least-squares problem. \\
iht & Iterative hard-thresholding solver for sparse inverse problems. \\
cgls & Conjugate-gradient least-squares solver with quadratic regularization. \\
seismic\_wavelet & Generate a normalized decaying sinusoidal seismic wavelet with edge tapering. \\
edge\_cosine\_taper & Apply a cosine/squared-sine taper to both ends of a vector. \\
\bottomrule
\end{longtable}

\section{Sampling operators}
This part of the API has two deliberately different meanings.

\subsection{Scalar/vector sampling}
For a vector $x\in\mathbb{R}^{n}$ and indices $i_1,\ldots,i_{n_s}$,
\[
   Sx = \begin{bmatrix}x_{i_1}&\cdots&x_{i_{n_s}}\end{bmatrix}^{T}.
\]
The adjoint scatters the values back to a length-$n$ vector.  If an index is
repeated, its adjoint contributions are summed.

\begin{minted}{julia}
n = 100
idx = [1, 5, 20, 75]
S = sampling_op(idx, n)

x = randn(n)
d = S * x                 # length(idx)
xa = S' * d               # length n
\end{minted}

\subsection{Whole-trace sampling for 2-D through N-D arrays}
For seismic-style arrays, dimension 1 is time.  Let
\[
   \mathrm{size}(X)=(n_t,n_2,n_3,\ldots,n_N),
   \qquad n_{\rm tr}=\prod_{j=2}^{N} n_j.
\]
Internally MiniOps uses
\[
   X_2 = \operatorname{reshape}(X,n_t,n_{\rm tr}).
\]
Then
\[
   SX = X_2[:,\mathrm{idx}],
\]
so the forward result always has size
\[
   n_t\times n_s, \qquad n_s=\mathrm{length}(\mathrm{idx}).
\]
The adjoint returns an array with the original full shape and sums repeated
trace indices.

\begin{minted}{julia}
# (time, x)
X2 = randn(1000, 120)
S2 = trace_sampling_op([1, 20, 75], size(X2))
D2 = S2 * X2                    # 1000 x 3
X2a = S2' * D2                  # 1000 x 120

# (time, x, y)
X3 = randn(1000, 30, 20)
S3 = trace_sampling_op([1, 31, 200], size(X3))
D3 = S3 * X3                    # 1000 x 3
X3a = S3' * D3                  # size(X3)

# (time, x, y, z, component)
X5 = randn(500, 10, 8, 4, 3)
S5 = trace_sampling_op([1, 17, 200], size(X5))
D5 = S5 * X5                    # 500 x 3
X5a = S5' * D5                  # size(X5)
\end{minted}

For spatial coordinates such as \texttt{(ix,iy)}, convert to the flattened
trace number with Julia's column-major convention:
\begin{minted}{julia}
j = LinearIndices((nx, ny))[ix, iy]
\end{minted}
The generic call \texttt{sampling\_op(idx,size(X))} is backward-compatible
and delegates to whole-trace sampling when \texttt{ndims(X) >= 2}.

\section{Operator examples}
\subsection{Convolution}
\begin{minted}{julia}
h = [1.0, -0.5, 0.25]
C = conv1d_op(h)
y = C * x
xa = C' * y
\end{minted}

For a gather with time along rows and traces along columns:
\begin{minted}{julia}
C = conv1d_cols_op(h)
Y = C * X
Xa = C' * Y
\end{minted}

\subsection{Scaling and diagonal weights}
\begin{minted}{julia}
S = scaling_op(0.5)
y = S * x

D = diag_op(w)
y = D * x
xa = D' * y
\end{minted}
Both operators use complex conjugation in their adjoints.

\subsection{FFT}
\begin{minted}{julia}
F = fft_op()                 # unitary by default
X = F * x
x2 = F' * X
\end{minted}
With the default normalization,
\[
  Fx = \frac{1}{\sqrt{N}}\operatorname{fft}(x),\qquad
  F^*X = \frac{1}{\sqrt{N}}\operatorname{bfft}(X).
\]
Therefore \texttt{F' * (F * x)} reproduces \texttt{x} to roundoff.

\subsection{Padding}
\begin{minted}{julia}
P = pad_op((100, 60), (20, 10), (20, 10))
Y = P * X
X2 = P' * Y
\end{minted}

\subsection{Blending}
\begin{minted}{julia}
source_times = [0.0, 0.37, 0.81]
B = blending_op(source_times, dt, nt)
b = B * D                       # D is nt x ns
Da = B' * b
\end{minted}

\section{NUDFT normalization}
The Taylor NUDFT operators use a \emph{unitary} Fourier convention.  In one
dimension,
\[
  Ux = \frac{1}{\sqrt{N}}\operatorname{fft}(x),\qquad
  U^*X = \sqrt{N}\operatorname{ifft}(X).
\]
Consequently, when all coordinate perturbations are zero,
\texttt{nudft1\_op(zeros(N),dx)} reduces to the unitary inverse transform
$U^*$, not Julia's raw \texttt{ifft}.  Thus
\[
  A X = \sqrt{N}\,\operatorname{ifft}(X)
  \qquad (\delta x=0).
\]
A direct comparison against \texttt{ifft(X)} will therefore differ by
$\sqrt{N}$.  This is a normalization convention, not an irregular-sampling
error.  The same convention is used by the multidimensional Taylor NUDFT.

\begin{minted}{julia}
N = 128
A = nudft1_op(zeros(N), dx)
X = randn(ComplexF64, N)

x1 = A * X
x2 = sqrt(N) .* ifft(X)
@show norm(x1 - x2) / norm(x2)
\end{minted}

\section{Radon operators}
\subsection{Parabolic Radon}
\begin{minted}{julia}
R = radon_tx_parab_op(dt, h, q; href=1000.0)
d = R * m                       # m: nt x nq
a = R' * d                      # a: nt x nq
\end{minted}

\subsection{Linear Radon by chirp factorization}
\begin{minted}{julia}
R = radon_tx_tp_chirp_op(dt, x0, dx, nx, p0, dp, np;
                         f1=2.0, f2=60.0)
d = R * m
ma = R' * d
\end{minted}

\section{Diagnostics}
Every newly written operator should first pass a linearity test and then an
adjoint test.
\begin{minted}{julia}
oklin, errlin = linearity_test(A, x1, x2)
okadj, erradj = adjoint_test(A, x, y)

sigma = opnorm_power(A, x; niter=30)
\end{minted}
For a correct double-precision implementation, the adjoint error is usually
close to machine precision, although interpolation, FFT conventions, and
single precision can justify a looser tolerance.

\section{Sparse solvers}
\subsection{ISTA}
\begin{minted}{julia}
sigma = opnorm_power(A, u0)
step = 0.9 / sigma^2
u = ista(A, y, u0, mu, step; niter=200)
\end{minted}
This minimizes an $\ell_1$-regularized least-squares objective using soft
thresholding.

\subsection{IHT}
\begin{minted}{julia}
u = iht(A, y, u0, tau, step; niter=200)
\end{minted}
IHT uses hard thresholding and is useful when explicit sparsity is desired.

\subsection{CGLS with quadratic regularization}
\begin{minted}{julia}
x = cgls(A, b, mu, x0; tol=1e-6, max_iter=500)
\end{minted}
The implemented objective is proportional to
\[
   \|Ax-b\|_2^2 + \mu\|x\|_2^2.
\]

\section{Source-file inventory and maintenance notes}
\begin{itemize}
  \item \texttt{core.jl}: operator type and shape helper.
  \item \texttt{algebra.jl}: adjoint, composition, scalar multiplication.
  \item \texttt{conv.jl}: 1-D and columnwise convolution.
  \item \texttt{sampling.jl}: scalar and whole-trace sampling.
  \item \texttt{blending.jl}: seismic blending.
  \item \texttt{scaling.jl}: scalar and diagonal operators.
  \item \texttt{fftops.jl}: FFT and padding operators.
  \item \texttt{nudftops.jl}: Taylor-series NUDFT operators.
  \item \texttt{radon\_parab.jl}: parabolic Radon operator.
  \item \texttt{radon\_chirp.jl}: chirp-factorized linear Radon operator.
  \item \texttt{diagnostics.jl}: operator validation and norm estimation.
  \item \texttt{solvers.jl}: ISTA, IHT, and CGLS.
  \item \texttt{extras.jl}: wavelet and taper utilities.
\end{itemize}

The file \texttt{radon\_linear.jl} is currently a legacy duplicate of the
parabolic implementation under older names. It is not included by
\texttt{MiniOps.jl} and therefore is not part of the active package API.
It should either remain clearly marked as legacy or be removed after any
needed comparison with older notebooks.

\section{Recommended validation before release}
\begin{enumerate}
  \item Run \texttt{Pkg.test()} under the target Julia version.
  \item Run adjoint tests with both real and complex data for operators that
        support complex arithmetic.
  \item Keep the 2-D through 5-D whole-trace sampling tests; they define the
        intended time-first sampling contract.
  \item When comparing NUDFT results with \texttt{ifft}, account explicitly
        for the unitary $\sqrt{N}$ normalization.
  \item Keep only active package files in release archives; editor swap
        files, \texttt{\_\_MACOSX}, and repository metadata are not needed.
\end{enumerate}

\end{document}
